\documentclass[../Main.tex]{subfiles}

\begin{document}
\section{Simple Methods for Approximating Integrals}
\subsection{Taylor Expansion of the Integrand}
We will consider Stieltjes integrals, which have the form:
\begin{equation*}
    I(x) = \int_{0}^{\infty} \frac{e^{-t}}{1 + xt} dt
\end{equation*}
for $x \geq 0$.

Our aim is to understand the behaviour of $I(x)$ for small $x$. We can Naively Taylor expand the integrand:
\begin{equation*}
    \frac{1}{1 + xt} =\sum_{n=0}^{\infty} (-1)^n x^n t^n
\end{equation*}
by unravelling a geometric series. However, our integral takes $t \to \infty$ and so this expression is not suitable. We could instead consider a finite sum:
\begin{align*}
    \sum_{n=0}^{N} (-xt)^n &= \frac{1 - (-xt)^{N+1}}{1 + xt} \\
    &= \frac{1}{1 + xt} - \underbrace{\frac{(-xt)^{N+1}}{1 + xt}}_{r_N}
\end{align*}
Substituting this into the integral:
\begin{align*}
    I(x) &= \int_{0}^{\infty} e^{-t} \left[\sum_{n=0}^{N} (-1)^n x^n t^n + r_N\right] dt \\
    &= \sum_{n=0}^{N} (-1)^n x^n \int_{0}^{\infty} t^n e^{-t} dt + R_n \\
    &= \sum_{n=0}^{N} (-1)^n x^n n! + R_N
\end{align*}
so we can define $\phi_n = x^n$, an asymptotic sequence. What remains now is to show that $R_N = o(x^N)$ as $x \to 0$. We can do this by considering a ratio:
\begin{align*}
    \frac{R_N}{x^N} &= \frac{x^{N+1}}{x^N} \int_{0}^{\infty} \frac{t^{N+1}e^{-t}}{1 + xt} dt \\
    &\leq x \int_{0}^{\infty} t^{N+1}e^{-t} dt \text{ since $1 + xt > 1$} \\
    &= x(N+1)!
\end{align*}
which tends to $0$ as $x \to 0$ because $N$ is fixed. Our final answer is:
\begin{equation*}
    I(x) \sim \sum_{n=0}^{\infty} (-1)^nn!~x^n \text{ as $n \to 0^+$.}
\end{equation*}
\subsection{Integration by Parts}
We can use the integration by parts formula:
\begin{equation*}
    \int_{a}^{b} f'(t) g(t) dt = \left[f(t)g(t)\right]_a^b - \int_{a}^{b} f(t)g'(t) dt 
\end{equation*}
many times:
\begin{align*}
    I &= \int_{a}^{b} f(t)g^{(n)}(t) dt \\
    &= \sum_{k=1}^{n} (-1)^{k-1} \left[f^{(k-1)}(t) g^{(n-k)}(t)\right]_a^b + (-1)^n \int_{a}^{b} f^{(n)}(t)g(t) dt
\end{align*}
This can generate an asymptotic approximation under certain circumstances.
\begin{example}
    Consider:
    \begin{equation*}
        F(x) = \int_{x}^{\infty} e^{-\frac12t^2} dt
    \end{equation*}
    We are interested in its behaviour as $x \to \infty$.

    We rewrite the integrand as $t \times \frac{e^{-t^2 / 2}}{t}$, and perform one integration by parts:
    \begin{equation*}
        F(x) = \left[-\frac1{t} e^{-t^2 / 2}\right]_x^\infty + \int_{x}^{\infty} \frac{e^{-t^2 / 2}}{t^2} dt.
    \end{equation*}
    A one-term asymptotic expansion is therefore $\frac{e^{-x^2 / 2}}{x}$, as long as we can bound the error:
    \begin{align*}
        \int_{x}^{\infty} \frac{e^{-t^2 / 2}}{t^2} dt &= \int_{x}^{\infty} \frac{te^{-t^2 / 2}}{t^3} dt \\
        &\leq \int_{x}^{\infty} \frac{te^{-t^2 / 2}}{x^3} dt \text{ because $t > x$} \\
        &= x^{-3} \int_{x}^{\infty} te^{-t^2 / 2} dt \\
        &= x^{-3} e^{-x^2 / 2}
    \end{align*}
    Therefore the error is $o\left(e^{-x^2 / 2} / x\right)$, and our approximation is valid.

    Continuing this process, we find the following asymptotic series:
    \begin{equation*}
        F(x) \sim \frac{e^{-x^2 / 2}}{x} \sum_{n=0}^{\infty} (-1)^n \frac{(2n-1)!!}{x^{2n}}
    \end{equation*}
\end{example}
\end{document}